\begin{proof}[Proof of Theorem~\ref{Theorem:2block}]

Consider a binomial relation obtained from two $\Lambda$-tableaux $T$ and $T'$ of size $3 \times d$ where $d >2$ whose contents are row-wise equal. By applying quadratic changes to the tableaux (changes involving only two columns) we will reduce the degree of this relation thus proving that the matching field ideal is quadratically generated. 

Given a $\Lambda$-tableau $T$, arrange the columns so that the first columns are those for which the matching field assigns the identity permutation and to the last column the matching field assigns the transposition $(12)$. Let $C$ denote the subtableau formed by the first columns and let $D$ denote the subtableau formed by the last columns. 

The tableaux $C$ and $D$ can each be put into semi-standard format. 
In other words, we can rearrange both $C$ and $D$ so that all rows are in weakly increasing order, and the columns of $C$ are strictly increasing, whereas the columns of $D$ are arranged so that the first and second entries are permuted from the diagonal order. 

Now given a binomial relation obtained from two $\Lambda$-tableaux $T$ and $T'$. We assume that $T$ (respectively $T'$) is organized as a pair of subtableaux $C$, $D$ (respectively $C'$, $D'$) satisfying the requirements described above. 
 If the first columns of $T$ and $T'$ are equal, then we can cancel them from the binomial relation and it is not a minimal generator. 
 We let $I$ and $I'$ denote the first columns of $T$ and $T'$, respectively. 
Otherwise by Lemma~\ref{lem:homo} the matching field relations are homogeneous with respect to the $\mathbb{Z}^4$-grading and so $\lvert I \cap I_1 \rvert = \lvert I' \cap I_1\rvert$.

\let\oldthecdrthm\thecdrthm
\setcounter{cdrthm}{0}
\def\thecdrthm{\arabic{cdrthm}}
\begin{enonce}[remark]{Case}\label{Case1}%\refstepcounter{cdrthm}
Suppose that $\lvert I \cap I_1 \rvert = \lvert I' \cap I_1\rvert = 3$. In this case, the columns $I$ and $I'$ could only differ in the second row. Suppose the entries of the second row of $I$ and $I'$ are $j $ and $j'$, respectively. We can also assume that $j <j'$. 
Then there must be a $j$ in the second row of the tableaux $D'$ since the contents are row-wise equal and $C'$ is in weakly increasing order. Then swap the positions of $j$ and $j'$ in the second row of $T'$ so that the first columns of $T$ and $T'$ now agree. Notice that we can exchange the position of $j$ with that of $j'$ since $j, j' \in I_1$ and $j'$ was originally in the second row of a column whose the first and third entries were in $I_2$. 
\end{enonce}

\begin{enonce}[remark]{Case}\label{Case2}%\refstepcounter{cdrthm}
Suppose that $| I \cap I_1 | = |I' \cap I_1| =2$. In this case, the columns $I$ and $I'$ may only differ in the second and third rows but not in the first. Assume the column $I$ is $i, j, r$ and the column $I'$ is $i, j', r'$. If $j <j'$ then just as above there must be a $j$ in the second row of $D'$. We have that $j <r'$ since $j$ is in the first block and $r'$ is in the second block, so we can swap the positions of $j $ and $j'$ in the second row of $T'$. 

Assume now that $j = j'$, without loss of generality we can suppose that $r<r'$. Then there is an $r$ in the third row of $D'$. Suppose that the column containing $r$ is $s, t, r$. Then we can swap the positions of $r$ and $r'$ since $t<r<r'$ and we can place $r$ in the last row. Now the two first terms are equal and hence, the binomial is not a minimal generator.
\end{enonce}

\begin{enonce}[remark]{Case}\label{Case3}%\refstepcounter{cdrthm} 
Suppose that $| I \cap I_1 | = |I' \cap I_1| = 1$. In this case, the tableaux $C$ and $C'$ are empty. Then the first two columns must be equal since $T = D$ and $T'= D'$ and they are both in (transposed) semi-standard form. 
\end{enonce}

\begin{enonce}[remark]{Case}\label{Case4}\setcounter{cdrthm}{4}
Suppose that $| I \cap I_1 | = |I' \cap I_1| =0$. In this case, the entries of $I$ and $I'$ can only differ in the first and third row. Suppose that the column $I$ is $r, s, t$ and that the column $I'$ is $r', s, t'$. 
Therefore $r, r'<s<t, t'$ and we can assume that $r' <r$. 
Then there is a column in $T'$ with $r$ in the first row and we can swap $r$ and $r'$ in the first row of $T'$. 
Thus we may assume that $r = r'$ and without loss of generality that $t < t'$. Then there must be a $t$ somewhere in the last row of $D'$ and we can again swap $t$ and $t'$ so that the columns are now equal. This completes the proof. \qedhere
\end{enonce}
\end{proof}
